\section{Quantitative preparation with varying bond dimensions}
\label{sec:ldp_rectangular}

Let $A_i^s:\C^{D_{i+1}}\to\C^{D_i}$ be the site matrices of a ring,
with $1\le D_i\le D$ and cyclic subscripts, and let $\sigma_i$ be
possibly singular density matrices on the bond spaces. Relative to the
inhomogeneous paragraph of~\cite{Malz2024Preparation}, the quantitative
hypothesis below acts on these rectangular matrices directly and is
sufficient for the claimed accuracy-dependent depth, though it is not
implied by the qualitative convergence condition in that paragraph.

\begin{lemma}[Rectangular blocking and zero padding]
    \label{thm:ldp_rectangular_blocking}
    \lean{MPSPreparation.blockSite_singleton,
        MPSChainTensor.rectangularEval,
        MPSChainTensor.rectangularTransfer,
        MPSChainTensor.rectangularTransfer_apply,
        MPSChainTensor.zeroPad_rectangularEval,
        VaryingBondChain.blockBondDim,
        VaryingBondChain.blockBondDim_site,
        VaryingBondChain.blockBondDim_succ,
        VaryingBondChain.blockBondDim_zero,
        VaryingBondChain.blockBondDim_length,
        VaryingBondChain.rectangularBlockSites,
        VaryingBondChain.rectangularBlockTensor,
        VaryingBondChain.zeroPad_rectangularBlockTensor}
    \leanok
    For a nonempty interval $[a,b)$, let
    $B^{\boldsymbol s}=A_a^{s_a}\cdots A_{b-1}^{s_{b-1}}$,
    a matrix of size $D_a\times D_b$. Let $J_i:\C^{D_i}\to\C^D$
    be the coordinate isometry, with $P_i=J_iJ_i^\dagger$.
    Then
    \begin{align}
        (J_aA_a^{s_a}J_{a+1}^\dagger)\cdots
        (J_{b-1}A_{b-1}^{s_{b-1}}J_b^\dagger)
         & =J_aB^{\boldsymbol s}J_b^\dagger,\notag \\
        F_{a:b}(X)
         & =\sum_{\boldsymbol s}B^{\boldsymbol s}X
        (B^{\boldsymbol s})^\dagger
        =F_a\circ\cdots\circ F_{b-1}(X).
    \end{align}
    The empty product on a bond space pads to $P_a$, whereas the
    empty product of ambient matrices is $I_D$.
\end{lemma}
\begin{proof}\leanok
    By induction on the interval length, each internal contraction uses
    $J_i^\dagger J_i=I_{D_i}$, and summing over the first physical index and
    then over the remaining word gives the transfer composition.
\end{proof}

For two adjacent sites, the following equality retains both physical
indices. The open virtual indices have dimensions $D_a$ and $D_{a+2}$;
the contracted middle bond has dimension $D_{a+1}$.
% TENKZ-RECTANGULAR-BLOCK-BEGIN
% Formula: B^(i,j)_(alpha,gamma) = sum_beta A_a^i_(alpha,beta) A_(a+1)^j_(beta,gamma).
% Boundary: alpha,gamma are virtual; i,j are physical. Only beta is contracted.
\begin{center}
    \tenkzkernel
    \begin{tenkzequation}
        \begin{tenkz}[rows={wire}, cols=4, bonds=none]
            \tn[at=(1,1), name=rectleft, skin=box,
                ports={180:virtual:$\alpha$, 0:virtual,
                        90:physical:$i$}]{A_a}
            \tn[at=(1,3), name=rectright, skin=box,
                ports={180:virtual, 0:virtual:$\gamma$,
                        90:physical:$j$}]{A_{a+1}}
            \tnwire{rectleft.0}{rectright.180}
        \end{tenkz}
        $=$
        \begin{tenkz}[rows={wire}, cols=3, bonds=none]
            \tn[at=(1,1), name=rectblock, skin=box, wide=2,
                ports={180:virtual:$\alpha$, 0:virtual:$\gamma$,
                        90@1:physical:$i$, 90@2:physical:$j$}]{B}
        \end{tenkz}
    \end{tenkzequation}
\end{center}
% TENKZ-RECTANGULAR-BLOCK-END

\begin{definition}[Corner reset references]
    \label{def:ldp_corner_reference}
    \lean{MPSPreparation.cornerProjection,
        MPSPreparation.cornerReferenceMap,
        MPSPreparation.cornerReferenceMap_apply,
        MPSPreparation.cornerProjection_full,
        MPSPreparation.cornerReferenceMap_full}
    \leanok
    For $\widehat\sigma_i=J_i\sigma_iJ_i^\dagger$, the padded reference
    for an interval $[a,b)$ is
    $Q_{a:b}(X)=\operatorname{Tr}(P_bX)\widehat\sigma_a$, which
    vanishes on inputs outside the outgoing bond corner.
\end{definition}

\begin{lemma}[Gram matrices and pairs on the bond spaces]
    \label{thm:ldp_rectangular_gram_pairs}
    \lean{MPSTensor.exists_norm_gram_transferMatrix_sub_le_of_entries,
        MPSTensor.zeroPad_eq_embeddedBlockState,
        MPSTensor.padPair_fixedPointPair,
        MPSPreparation.posSemidef_cornerProjection,
        MPSPreparation.cornerProjection_mul_self,
        MPSPreparation.sqrt_transpose_kronecker_corner,
        MPSPreparation.exists_norm_polarPos_sub_le_sqrt_transferMatrix_of_square,
        MPSPreparation.exists_norm_polarPos_sub_le_sqrt_cornerReferenceMap,
        MPSPreparation.ofPhysicalMatrix_sqrt_transpose_kronecker_corner,
        MPSPreparation.mixedMapLM_sqrt_transpose_kronecker_corner,
        MPSPreparation.exists_norm_transferMatrix_zeroPad_sub_cornerReferenceMap_le}
    \leanok
    Let every map between the bond matrix spaces carry its Hilbert--Schmidt
    induced operator norm, namely the $L^2$ operator norm of its
    matrix in the matrix-unit bases. Padding changes transfer-matrix
    errors by at most a constant depending only on $D$.
    The Gram reference and its positive square root are
    \begin{align}
        G_{a:b} & =\widehat\sigma_a^{\mathsf T}\otimes P_b,
                & \sqrt{G_{a:b}}                            & =(\sqrt{\widehat\sigma_a})^{\mathsf T}\otimes P_b.
    \end{align}
    Consequently the positive-polar error is bounded by a
    dimension-dependent constant times
    $\|F_{a:b}-\operatorname{Tr}(\cdot)\sigma_a\|^{1/2}$.
    Also $\sqrt{\widehat\sigma_i}=J_i\sqrt{\sigma_i}J_i^\dagger$,
    so the pair obtained from the padded reference is the
    padding of the pair on the bond space.
\end{lemma}
\begin{proof}\leanok
    \uses{thm:ldp_rectangular_blocking,thm:ldp_sqrt_holder}
    Extending the matrix-unit coefficient matrix by zero and permuting its
    indices into the transfer convention gives a linear map between
    endpoint spaces of dimension at most $D$; since there are only finitely
    many such dimension pairs, these linear maps have a common
    operator-norm bound. Reshuffling the corner reset gives $G_{a:b}$,
    whose displayed square root is positive and squares to it, and the
    global square-root H\"older estimate then bounds the positive-polar
    error. Since conjugation by a coordinate isometry preserves positivity
    and squares, the square-root and pair identities follow.
\end{proof}

\begin{lemma}[Ordered corner references and normalized error]
    \label{thm:ldp_rectangular_error}
    \lean{MPSPreparation.transferMatrix_cornerReferenceMap_eq_mul_diagonal,
        MPSPreparation.norm_transferMatrix_cornerReferenceMap_le,
        MPSPreparation.transferMatrix_cornerReferenceMap_mul,
        MPSPreparation.prod_range_transferMatrix_cornerReferenceMap,
        MPSPreparation.trace_transferMatrix_cornerReferenceMap,
        MPSPreparation.exists_uniform_reference_bounds,
        MPSPreparation.exists_norm_trace_prod_sub_one_le_of_reference_bounds,
        MPSPreparation.exists_norm_trace_prod_polarPos_sub_one_le_corner_reference,
        MPSPreparation.exists_abs_norm_chainState_sq_sub_one_le_corner_reference,
        MPSPreparation.exists_one_sub_norm_inner_chainPosState_le_varying_reference}
    \leanok
    For consecutive blocks, $Q_{a:b}Q_{b:c}=Q_{a:c}$, and the matrix
    trace of a whole-ring reference is one. Every contiguous reference
    product has norm bounded in terms of $D$ alone.
    If every blocked transfer and the whole-ring transfer have error
    at most $\delta$ from their corner references, the normalized
    positive-part state has pair error at most $C_DM\sqrt\delta$.
    The pair leaving a block uses the density at the next block's
    incoming bond.
\end{lemma}
\begin{proof}\leanok
    \uses{thm:ldp_rectangular_gram_pairs,thm:ldp_varying_reference_products,thm:ldp_varying_reference_pairs,thm:ldp_normalization_triangle}
    Because $\operatorname{Tr}(P_b\widehat\sigma_b)=1$, composition removes
    the middle bond and leaves a contracting outgoing-corner mask. The
    ordered-product perturbation estimate applied to the mixed transfers,
    followed by the trace, bounds the pair overlap, while the trace of the
    whole-ring transfer error controls normalization. The inverse-norm
    estimate then gives the stated normalized error, including the regime
    where the trivial error bound one is sharper.
\end{proof}

\begin{lemma}[Physical logarithmic depth from rectangular mixing]
    \label{thm:ldp_rectangular_preparation}
    \lean{mul_exp_neg_mul_le_of_div_log_le,
        MPSPreparation.exists_isPreparedInDepth_le_log_eventually_of_rectangular_mixing}
    \leanok
    Fix $d>0$ and $D$, and let the preceding data vary with the ring length
    $N$. Suppose $K,r>0$ are independent of $N$ and every nonempty
    block $[a,b)$ of every partition of the ordered ring satisfies
    \begin{align}
        \|F_{a:b}-\operatorname{Tr}(\cdot)\sigma_a\|
        \le K e^{-r(b-a)}.
    \end{align}
    There are $N_0$ and $C$, chosen before $N$ and $\varepsilon$, such
    that for positive ring lengths $N\ge N_0$ and $0<\varepsilon\le1$, the normalized ring
    state admits a unit-vector approximation with error at most
    $\varepsilon$ prepared by a ring circuit of depth at most
    $C\log(N/\varepsilon)$.
\end{lemma}
\begin{proof}\leanok
    \uses{thm:ldp_rectangular_error,thm:ldp_inhomogeneous_uniform_rate_from}
    Whole-ring mixing gives
    $|\|\phi_N\|^2-1|\le C_DKe^{-rN}$, hence nonvanishing beyond
    a fixed cutoff. Partitioning into blocks with
    $q\le\ell_j\le2q$ bounds the pair error by
    $C'_D\sqrt K\,N e^{-rq/2}$, and the uniform-rate preparation theorem
    then selects a logarithmic block length, falling back to exact
    preparation if that length exceeds $N$.
\end{proof}

Equation (S39) of~\cite{Malz2024Preparation} concerns open-boundary
states in its numerical study. Dimension-one bonds are allowed in the
preceding theorems, but their physical circuits use ring adjacency.
An open-line gate conclusion would additionally require a construction
that avoids the wrap-around edge.
